\section{Introduction}


Recent years have seen a proliferation of new ideas and techniques in stable homotopy theory, particularly in  equivariant techniques and algebraic K-theory.  At the intersection of these two areas is \emph{genuine equivariant algebraic $K$-theory}, which studies invariants of rings (and ring spectra) with an action by a compact Lie group $G$. Equivariant algebraic $K$-theory is related to special values of Artin $L$-functions, and plays a central role in Malkiewich and Merling's recent work on a refinement of the stable parametrized $h$-cobordism theorem \cite{ElmantoZhang,MM1,MM2}. It has also been effectively utilized by Vogeli to perform new, non-equivariant computations in the algebraic $K$-theory of group algebras \cite{Vogeli}. 
Unfortunately, computations in non-equivariant algebraic K-theory are difficult, and are often much more difficult in the equivariant context.

One very successful approach to computation in algebraic $K$-theory is trace methods, in which one approximates the algebraic $K$-theory of a ring $R$ using the \emph{cyclotomic trace} $K(R)\to \TC(R)$, which is often a good approximation.  Here, $\TC(R)$ denotes \emph{topological cyclic homology}, which is derived from \emph{topological Hochschild homology}, denoted $\THH(R)$, using the fact that $\THH(R)$ is a cyclotomic spectrum. The \emph{Dennis trace} gives a map $K(R)\to \THH(R)$, which plays a crucial role in the construction of the cyclotomic trace. 

A natural question to ask is whether the tools from trace methods can be developed to aid in computations of the equivariant algebraic $K$-theory of a $G$-ring $R$. In \cite{CGK25}, some of the authors took a step in this direction, introducing equivariant THH, denoted $\ETHH$, which receives an equivariant refinement of the Dennis trace, $K_G(R)^G\to\ETHH(R)^G$.  Here, $K_G(R)$ is the equivariant algebraic $K$-theory $G$-spectrum, in the sense of Merling \cite{Merling}. 

Recent work of Hilman--Ramzi \cite{HilmanRamzi} develops a notion of equivariant THH in a slightly different setting; upcoming work of the second named author will show equivalence of the definitions for $G$-ring spectra. Hilman--Ramzi also develop an equivariant Dennis trace from the equivariant algebraic $K$-theory of Barwick et al.\ \cite{Bar17, BGS20} to their version of equivariant THH. The relationship between the Dennis trace maps of \cite{HilmanRamzi} and of \cite{CGK25} is unclear.

Equivariant $\THH$ comes equipped with a circle action, and forthcoming work of the second named author will use the circle action to describe the sense in which $\ETHH(R)$ is cyclotomic and use this structure to define an equivariant notion of $\TC$ \cite{Got}.

While we now have an equivariant refinement of the Dennis trace, the effectiveness of a trace methods approach depends heavily on the computability of the homotopy groups of $\ETHH(R)$.  The goal of the present paper is to provide computations of homotopy groups $\ETHH(R)$ for several fundamental ring $G$-spectra.

While we now have an equivariant refinement of the Dennis trace, the effectiveness of a trace methods approach depends heavily on the availability of fundamental computations of both $\ETHH(R)$ and of equivariant $K$-theory from which to build upon.  While some computations have been carried out,  for instance \cite{AGHKK22,ChanVogeli,CW25,Wis26}, the literature of explicit computations remains somewhat sparse. With this in mind, the goal of the present paper is to provide computations of the homotopy groups of $\ETHH(R)$ for several fundamental ring $G$-spectra.

\subsection{Equivariant B\"okstedt Periodicity}

The first computations of THH were done by B\"okstedt, who computed $\THH(H\mathbb{Z})$ and $\THH(H\F_p)$ for all primes $p$. These computations play a fundamental role in the theory, and continue to be important building blocks for new computations \cite{KrauseNikolaus:BokstedtPeriodicity,BhattMorrowScholze}.  In the $C_p$-equivariant setting, the correct analogue of $\mathbb{F}_p$ is the \emph{constant Green functor} $\underline{\F}_p$, and our first main theorem computes the homotopy groups of the $C_p$-geometric fixed points of $\ETHH(\HFp)$.

We write $\Gamma$ and $\Lambda$ for divided power and exterior algebras over $\F _p$, respectively.

\begin{introtheorem}[{\cref{lemma: injection from gen to geo,thm:Bokstedt}}]\label{thm:introBokstedt}
    Let $p$ be an odd prime. The graded homotopy ring $\pi^{C_p}_*(\ETHH(\HFp))$ is isomorphic to the subring of $\Gamma[a]\otimes\F_p[b,c]\otimes \Lambda[d,e]$ generated by $b$ and all monomials divisible by either $e$ or a divided power of $a$. The degrees are $|a|=|b|=|c|=2$, $|d|=1$, and $|e|=3$.
\end{introtheorem} 

In comparison, B\"okstedt's foundational result showed that $\pi_*\THH(H\F _p) \cong \F _p [b]$, where $b$ is a polynomial generator in degree 2. This result is known as B\"okstedt periodicity. Our theorem demonstrates that equivariant B\"okstedt periodicity is far more complicated. 

We prove \cref{thm:introBokstedt} using a careful analysis of the Tate square,
\[
    \begin{tikzcd}
        \ETHH(\HFp)^{C_p}\ar[r,"i"] \ar[d] & \ETHH(\HFp)\geop \ar[d]\\
        \ETHH(\HFp)^{hC_p} \ar[r] & \ETHH(\HFp)^{tC_p},
    \end{tikzcd}
\]
which is a pullback of commutative ring spectra. The geometric fixed points are interesting in their own right.


\begin{introtheorem}[{\cref{prop-THH-geomfp,cor: homotopy of geo of ETHH hfp}}]\label{introthm: geo of ETHH hfp}
    Let $p$ be an odd prime.  There is an equivalence of $\mathbb{E}_1$-$H\F_p$-algebras 
    \[
        (\ethh(\HFp))\geop \simeq H\F_p \wedge (\Omega S^3)_+ \wedge S^1_+ \wedge \C P^\infty _+ \wedge (LS^3)_+ 
    \]
    where $LS^3$ denotes the free loop space. Consequently, there is an isomorphism of graded rings
    \[
        \pi_*\ETHH(\HFp)\geop\cong \Gamma[a]\otimes\F_p[b,c]\otimes \Lambda[d,e]
    \]
    where $a$, $b$, $c$, $d$, and $e$ have the same degrees as in \cref{thm:introBokstedt}.
\end{introtheorem}

The proof of \cref{introthm: geo of ETHH hfp} utilizes the fact that $\ETHH$ commutes with the geometric fixed points construction (see \cref{cor-geomfp-ETHH}). In particular, this reduces the computation to the non-equivariant computation of $\THH(\HFp\geop)$. From here, we prove an equivalence of $\mathbb{E}_2$-$H\F_p$-algebras $\HFp\geop\simeq H\F_p\wedge S^1_+\wedge \Omega S^3_+$ and use the monoidal properties of THH to prove the theorem.

Returning to the proof of \cref{thm:introBokstedt}, we see that on homotopy, the fixed points are a subring of the geometric fixed points. To fully pin down the subring, the final, key part of the analysis is understanding how the right vertical map acts on the generators of the homotopy of $\ETHH(\HFp)\geop$.




When $p=2$, the $C_2$-homotopy groups of $\ETHH(H\underline{\F}_2)$ are computed in \cite{AGHKK22} using the fact that $H\underline{\F}_2$ is a sufficiently multiplicative Thom spectrum \cite{BW18,HW20}. The computation of the ring structures on the homotopy groups of $\ETHH(H\underline{\F}_2)\geo[C_2]$ and $\ETHH(H\underline{\F}_2)^{C_2}$ is more difficult and requires a different method than the one we apply here in the odd prime case.  This computation will appear in forthcoming work of the first-named author and Chase Vogeli. 

As a future direction, we think it would be interesting to understand the $RO(C_p)$-graded homotopy ring of $\ETHH(\HFp)$.  Oftentimes, this ring has somewhat better formal properties than the integer graded homotopy ring, and could shed some light on the somewhat unorthodox answer obtained in \cref{thm:introBokstedt}. Additionally, an understanding of the $RO(C_p)$-graded homotopy groups would give access to norm structures on homotopy groups which could prove useful for later computation related to equivariant topological cyclic homology.




\subsection{Equivariant complex cobordism}

One of the most important ring spectra is the complex cobordism spectrum $MU$, which plays a fundamental role in computations of stable homotopy groups. The homotopy $S^1$-fixed points of $\THH$ of some quotients of $MU$ have been used to resolve the telescope conjecture \cite{BHLS23} and to study redshift \cite{HW22}. A full computation of $\TC(MU)$ could be used to compute information about $\TC$ of the sphere spectrum, an invariant with applications to geometric topology \cite{Hesselholt:Whitehead,Rognes:Cobordism}. The equivariant refinement of the complex cobordism spectrum, denoted $MU_G$, has been studied extensively; see, for instance, \cite{GreenleesMay:Localization,Hausmann}. Note that we are working with tom Dieck's homotopical equivariant cobordism (a Thom spectrum carrying the universal $G$-equivariant complex orientation \cite{CGK02}) as opposed to geometric cobordism (which encodes information about cobordism classes of $G$-manifolds).

The non-equivariant computation of $\THH(MU)$ can be done using the fact that $MU$ is a Thom spectrum.  In particular, \cite{BCS10} gives an equivalence
\[
    \THH(MU)\simeq MU\wedge SU_+
\]
where $SU$ denotes the infinite special unitary group. From here, the homotopy groups can be computed using the Atiyah--Hirzebruch spectral sequence.

In this paper, we upgrade the compatibility of equivariant factorization homology with Thom spectra proven in \cite{HHKWZ24} to include stronger multiplicative properties. We use this to compute $\ETHH(MU_G)$ as a $G$-$\mathbb{E}_\infty$ ring spectrum.
\begin{introtheorem}[{\cref{cor:identification-of-ETHH-of-MU_G}}]\label{introthm:identification-of-ETHH-of-MU_G}
    Let $G$ be a finite group. There is an equivalence
    \[
        \mathrm{Res}_{G}^{G \times S^1} \mathrm{ETHH}(MU_G) \simeq MU_G \wedge \Sigma^\infty _+ SU_G
    \]
    of $G$-$\mathbb{E}_\infty$ ring spectra, where $SU_G$ denotes an appropriately equivariant infinite special unitary group.
\end{introtheorem}

When $G$ is abelian, we use this to give a computation of the $G$-equivariant homotopy groups of $\mathrm{ETHH}(MU_G)$. 

\begin{introtheorem}[{\cref{thm:MU_G-homology-of-SU_G-for-G-abelian}}]\label{introthm: homotopy of ETHH(MU_G)}
    Let $G$ be compact, abelian Lie group and $H \subset G$. Then we have
    \[
        \pi^H_*(MU_G \wedge (SU_G)_+) \cong (MU_H)_* \otimes_\Z \Lambda(\lambda_1,\lambda_2,...)
    \]
    where $\lambda_i$ is an exterior algebra generator in degree $2i+1$. In light of \cref{introthm:identification-of-ETHH-of-MU_G}, this computes $\pi^H_*(\ETHH(MU_G))$ when $G$ is finite abelian.
\end{introtheorem}

The proof of \cref{introthm: homotopy of ETHH(MU_G)} relies on an explicit $G$-CW structure for $SU_G$ which we construct in \cref{section: MUG}.  This cell structure is a $G$-equivariant refinement of a CW-structure on $SU$ due to Yokota \cite{Yok56}.

When $G=C_2$, we also consider the Real analogue of complex cobordism, denoted $MU_{\R}$.  This spectrum is crucial to geometric applications of equivariant homotopy. For instance, it plays a central role in Hill, Hopkins, and Ravenel's work on the Kervaire invariant one problem \cite{HHR}. We are able to give a complete description of $\ETHH(MU_{\R})$, as well as its $RO(C_2)$-graded homotopy ring.

\begin{introtheorem}[{\cref{cor-ETHH-MUR,cor: homotopy of ETHHMUR}}]\label{introthm-ETHH-MUR}
    There is an equivalence of $C_2$-$\mathbb{E}_\infty$ ring spectra
    \[
        \mathrm{Res}_{C_2}^{C_2 \times S^1} \mathrm{ETHH}(MU_\mathbb{R}) \simeq MU_\mathbb{R} \wedge \Sigma^\infty _+ B(BU_\mathbb{R})
    \]
    where $BU_{\mathbb{R}}$ denotes the classifying space of stable Real vector bundles.  Consequently, there is an isomorphism 
    \[ 
        \underline{\pi}_\bigstar \Res_{C_2}^{C_2 \times S^1} \ETHH(MU_\R) \cong \Lambda_{(\underline{MU_\R})_\bigstar} (\overline{\lambda}_n | n \geq 1)
    \]
    of $RO(C_2)$-graded Green functors, where $\overline{\lambda}_n$ has degree $n \rho+1$.
\end{introtheorem}


This computation uses the multiplicative $G$-equivariant bar spectral sequence (\cref{thm-ebar,thm-ebar-multiplicative}), which may be of independent interest.

\subsection{Outline}

The remainder of the paper is organized as follows. Background material on ETHH, equivariant Thom spectra, and factorization homology can be found \cref{sec: prelims}.  In \cref{sec:barss}, we construct a multiplicative $G$-equivariant bar spectral sequence which we need for computing $\ETHH(MU_\R)$.  The proofs of \cref{introthm-ETHH-MUR,introthm:identification-of-ETHH-of-MU_G,introthm: homotopy of ETHH(MU_G)} are in \cref{sec: ETHH of Thom}.  Finally, we prove \cref{introthm: geo of ETHH hfp,thm:introBokstedt} in \cref{sec: Bokstedt}.


\subsection{Acknowledgments}
We would like to thank Ben Antieau, Teena Gerhardt, Mike Hill, Liam Keenan, Tyler Lawson, Maximilien P\'eroux, J.D. Quigley, and Chase Vogeli for helpful conversations related to this work.  DC and MG were partially supported by NSF grant DMS-2135960. IK is grateful for funding from the NWO-XL grant ``Symmetry on the interface of topology and higher algebra", which supported her travel for collaboration on this project. %\davidnote{add grant stuff.}